\section{Results}

\subsection{A single-nucleus atlas of the LUAD invasion sequence}

Paired single-nucleus RNA sequencing and Visium spatial transcriptomics were obtained from
23 patients spanning the five stages of the sequence, with the spatial arm covering 56
sections from 25 patients and sharing 23 patients with the single-cell arm
(Fig.~\ref{fig:1}a). Quality control retained 399,579 nuclei. The atlas resolves into six
lineages (epithelial, fibroblast, endothelial, myeloid, T/NK, B/plasma;
Fig.~\ref{fig:1}a) and, on per-lineage sub-clustering, into 39 second-level (L2) subtypes
(Supplementary Table~S1; Fig.~\ref{fig:1}f).

Two properties of the reference annotation affected the analyses that follow and are
recorded here.

\emph{Endothelial and myeloid panels were rebuilt.} The published marker table supplies six
canonical genes per endothelial subtype, which is insufficient to discriminate between
closely related endothelial states. Under that panel the \emph{Artery} subtype was assigned
to 16 clusters across the cohort; under an expanded panel of up to 20 genes per subtype
across nine endothelial types it was assigned to four, with the remainder reassigned to
\emph{Lymphatic}, \emph{Vein} and \emph{Capillary}. The endothelial and myeloid panels were
therefore rebuilt before any downstream analysis. Analyses that depend on endothelial
subtype proportions inherit the corrected panels.

\emph{The L2 nomenclature is not a set of unique identifiers.} The 39 names map onto 67
distinct lineage--subtype combinations. The clearest case is \emph{Basophil/Mast 1}, which
the reference assigns predominantly to epithelial cells (6,360 cells) and only marginally
to myeloid cells (8 cells), so the name misrepresents the content. Subtype composition is
accordingly reported at the lineage level (Fig.~\ref{fig:1}e).

A second difficulty concerns the coverage rather than the nomenclature of the reference. The
39-subtype reference contains no exhaustion programme, and exhausted T cells were
consequently absorbed into the nearest available label. One cluster carries LAG3, TIGIT,
HAVCR2, ENTPD1, GZMB and PRF1 while lacking CCR7, SELL, TCF7 and LEF1, an exhausted rather
than a naive phenotype, yet the reference assigns it to \emph{CD8+ Naive T}. Deconvolution
cannot recover a state absent from the reference, so the T-cell composition reported here is
a composition in the reference's vocabulary. Functional exhaustion was therefore measured
independently, by projecting an exhaustion signature
\citep{thommen2018,wherry2015exhaustion} onto the spots.

One AAH sample (P24) carried a plasma-cell programme confined to a single patient and not
tracking disease stage. A patient-level rank of B/plasma content, comparison against a
matched negative-control gene set, and an independent readout from a separate region of the
same block each identified it as ambient RNA rather than a tumour-associated infiltrate. The
sample was retained with an annotation, and all conclusions that depend on B/plasma content
are reported both with and without it. The source study does not describe ambient-RNA
correction, and P24 is one of its representative samples.

\begin{figure*}[!ht]\centering
  \includegraphics[width=\textwidth]{Fig1}
  \caption{\textbf{A single-nucleus atlas of the LUAD invasion sequence.}
(\textbf{a}) UMAP of 399,579 nuclei coloured by lineage.
(\textbf{b}) Lineage composition across the cohort.
(\textbf{c}) The same embedding coloured by disease stage.
(\textbf{d}) The embedding split by stage, retaining lineage colours.
(\textbf{e}) Lineage composition per stage.
(\textbf{f}) Composition of the 39 second-level subtypes per stage.
Endothelial and myeloid marker panels were rebuilt before composition analysis (Methods).
The 39 subtype names are not unique lineage--subtype identifiers (67 combinations);
compositions are reported at the lineage level in (\textbf{e}).}\label{fig:1}
\end{figure*}


\subsection{Spatial deconvolution resolves seven recurring niche domains}

BANKSY applied per section yielded 434 spatial domains, which collapsed to seven archetypes
($K^{*}=7$; cross-seed stability 0.945) across the 56 sections (Fig.~\ref{fig:2}d).
Stability falls to 0.814 at $K=8$, so the domains cannot be subdivided further without
becoming seed-dependent. We therefore treat $K^{*}=7$ as the resolution limit of this
cohort. A purely descriptive optimum exists at $K=6$ (stability 0.966) and was not adopted.

Identities were assigned from each archetype's RCTD composition (Fig.~\ref{fig:2}a,c) and
its marker genes (Fig.~\ref{fig:2}g). Six of the seven are ordinary lung structures: airway
epithelium (D1), peribronchovascular interstitium carrying an inflammatory CAF programme
(D2) \citep{ohllund2017,elyada2019}, alveolar wall (D4), AT2 epithelium (D5), vascular and
airway smooth muscle (D6), and lymphocyte aggregates (D7). Only D3 is restricted to
invasive disease. On the stage axis (Fig.~\ref{fig:2}d), D3 accounts for 17\% of spots in
IAC and \textbf{0\% in every precursor stage}, whereas D7 is present in all five stages.
The presence of lymphoid aggregates before invasion argues against the assumption that
immune aggregation accompanies malignancy.

Two caveats attach to the domain labels. The published marker table for these domains is
itself depth-confounded (see below), and the identity of D3 was revised twice during the
study, from an immune-infiltrate reading to a hypoxic/invasive reading and then to its
current ECM/interstitial reading; the earlier names should not be used. The domain
definition is also method-sensitive: BANKSY, RCTD composition and marker-based clustering
return 7, 6 and 5 domains respectively, with an adjusted Rand index of only 0.225--0.247
against an RCTD-aligned partition. Domain boundaries are therefore quantitative but not
unique. The composition underlying each label is given in Table~\ref{tab:domains}.

\begin{figure*}[!ht]\centering
  \includegraphics[width=\textwidth]{Fig2}
  \caption{\textbf{Spatial deconvolution resolves seven niche-domain archetypes.}
(\textbf{a}) RCTD lineage composition of each domain archetype (56 sections).
(\textbf{b}) Spatial matrix: one representative section per stage (rows) against
H\&E, RCTD at the 6-lineage and 39-subtype level, the four spatial programmes
(dominant programme per spot) and the niche-domain archetype (columns). Scale bars are
calibrated per section from the nominal 100-$\mu$m spot spacing.
(\textbf{c}) RCTD subtype composition of each domain archetype.
(\textbf{d}) Domain composition per stage; D3 is confined to invasive disease
(17\% of spots, 0\% in all precursor stages) whereas the lymphoid domain D7 occurs in all
five stages.
(\textbf{e}) Relative make-up and absolute level of the four spatial programmes per stage.
\textbf{This panel is a negative control}: before depth matching the programme scores
essentially mirror sequencing depth, and after matching into depth deciles the stage
ordering is identical across all ten deciles but does not follow disease progression. The
fourth programme is the lipid-associated macrophage (LAM) signature \citep{cheng2021trem2lam},
which we report alongside three exhaustion axes; it is not an exhaustion signature in the
T-cell sense.
(\textbf{f}) Univariable Cox regression of each domain's marker signature on overall and
disease-free survival (TCGA-LUAD); p\,$<$\,0.05 in red. The D3 row uses the unadjusted
signature (see main text).
(\textbf{g}) Marker-gene dot plot for the seven domains.}\label{fig:2}
\end{figure*}


\begin{table*}[!ht]
\centering
\begin{tabularx}{\textwidth}{@{}>{\bfseries}l >{\raggedright\arraybackslash}p{86pt} >{\raggedright\arraybackslash}p{152pt} >{\raggedleft\arraybackslash}p{50pt} >{\itshape\raggedright\arraybackslash}X@{}}
\toprule
\tabheadfont Domain & \tabheadfont Label & \tabheadfont Basis (RCTD enrichment or depth-balanced markers) & \tabheadfont Share & \tabheadfont Note \\
\midrule
D1 & Airway epithelium & Ciliated $\times$6.3; Goblet $\times$5.6 & 5\% & \\
D2 & iCAF stroma & IL6, CXCL8, CCL2, LIF, HAS1, HAS2 & 11\% & peribronchovascular \\
D3 & ECM / interstitial & COL1A1, COL3A1, FN1, CTHRC1 & 17\% (IAC) & no precursor stage; two label revisions \\
D4 & Alveolar wall & Capillary $\times$2.5; AT1 $\times$1.9 & 45\% & 77\% of the Normal section \\
D5 & AT2 epithelium & AT2 weight 0.480 ($\times$2.0) & 19\% & highest copy-number score \\
D6 & Vascular / airway wall & VSMC $\times$5.1; Artery $\times$2.3 & 8\% & \\
D7 & Lymphoid aggregate & B $\times$9.1; NKT $\times$4.7 & 4\% & present in all five stages \\
\bottomrule
\end{tabularx}
\caption{\textbf{Identity of the seven niche-domain archetypes.} Enrichment is the RCTD
composition of a domain relative to the rest of the cohort, as a fold change; the share is
the fraction of all spots assigned to that domain. The D3 label was revised twice during
the study (see main text) and the earlier names should not be used. No label here implies a
malignancy call.}
\label{tab:domains}
\end{table*}


\subsection{Sequencing depth is a proxy for tissue density}

The seven domains differ in sequencing depth by a factor of 1.00 to 7.95. Depth matching
removes a large and predictable fraction of each domain's marker genes: across the seven
domains, the fraction of a domain's top-150 up-genes surviving per-slide depth matching
falls monotonically with depth (Spearman $\rho=-0.89$; Fig.~\ref{fig:3}f), with two
independent estimators (decile matching and per-gene regression on $\log_{10}$(nUMI))
agreeing almost gene-for-gene. D3 is the sole outlier; it is both the deepest domain and the
only one restricted to invasive disease, and 16 of its 150 markers survive matching.

This resembles a standard technical confound, and treating it as one would be an error. The
same quantity is independently visible in the images: the PLIP-derived neoplastic score
correlates with the fraction of a spot occupied by tissue at $\rho=0.754$ across four
sections of one patient. Deep spots are therefore deep because they are dense, and dense
because they contain tumour. Removing depth as a covariate removes part of the biological
signal; the operation is closer to over-adjustment than to deconfounding, and neither the
unadjusted nor the adjusted estimate is correct on its own.

This coupling is the binding constraint of the study and recurs in four places. In the
spatial copy-number arm it determines which comparisons survive. In the prognosis arm it is
the reason one domain's apparent prognostic effect shrinks when its signature is replaced by
a depth-balanced version. In the developmental analysis it accounts for a trajectory running
opposite to that of the source study. It also prevents interpretation of the
stage-dependence of the four spatial programmes. Before matching, their scores are almost a
mirror of depth (per-spot $r=-0.30$ to $-0.69$ for the T-cell programme; sections differ
2.4-fold in median depth). After matching into depth deciles the stage ordering is identical
across all ten deciles but does not follow disease progression: AAH and AIS are both
precursor stages yet occupy opposite ends of the ordering. That comparison is consequently
reported as a negative control (Fig.~\ref{fig:2}e).

Because the revised identity of D3 is an ECM/interstitial programme, that programme was
tested directly rather than accepted from the domain analysis. Within each patient,
fibroblasts were scored against the cohort fibroblast mean and standard deviation per gene,
and the mean of a 22-gene ECM set was compared between invasive and precursor fibroblasts.
The programme is higher in invasive disease in \textbf{20 of the 22 patients with both
stages} (median within-patient difference $+0.402$; 300 random gene sets of matched size
drawn from expressed genes give $P=0.0033$ on both the count and the median). Eight of the
22 genes were added by us rather than drawn from the depth-balanced D3 ranking; restricting
the panel to the 14 data-derived genes gives 20 of 22 ($+0.347$, $P=0.0033$), restricting it
to the eight hand-added genes also gives 20 of 22 ($+0.450$, $P=0.0033$), and removing any
single gene leaves the count at 20 of 22 in all 22 leave-one-out variants. The ECM programme
is therefore not an artefact of gene selection, although the result describes a tissue-level
programme in fibroblasts and does not identify which ECM component is responsible.

\subsection{Copy-number inference: three failures and one surviving signal}

Malignant-cell identification is where the analytical failures of this study concentrate,
and they are reported in full because they constrain the interpretation of the spatial
results.

\paragraph{Single-cell inference.} Across six tool families, comprising CopyKAT
\citep{gao2021copykat} in default and floored configurations, two anchor-based arms, inferCNV
v1 and v2 \citep{patel2014infercnv}, and a supervised expression-based classifier
\citep{scmalignantfinder2025}, no configuration produced a usable malignant/non-malignant
separator, and each failure was traced to a specific cause. The clearest was an instability
check: the same set of 3,890 normal cells, used as the reference in three separate runs,
produced aneuploid calls of 0\%, 13.6\% and 49.8\%. The most misleading was an AUROC of
exactly 0.5000 from inferCNV v1, where all 19,748 cells carried a single identical score
value because the per-cell statistic reduced to the constant reference baseline; the ROC was
a degeneracy rather than evidence of an absent difference. A second version produced AUROC
0.652 but failed its own control, the normal cells scoring 0.0826 against a threshold of
0.05.

\paragraph{Cohort-level spatial copy number.} All 25 cases ran successfully at the Visium
level. The cohort-level claim that invasive sections carry more copy-number signal shrank by
71\% after depth balancing, and the matched comparison was not significant ($p=0.092$,
confidence interval including zero). A random subsample reproduced the pre-matching result
index by index, confirming that the original effect was driven by sequencing depth rather
than by sample size. No cohort-level copy-number difference is therefore reported.

\paragraph{One comparison survives by design.} Comparing each domain against the remaining
spots of the same section removes depth by construction, since every spot in the comparison
shares one cDNA library (Fig.~\ref{fig:3}c). Under this design the copy-number score is
elevated in D5 (AT2; median $\Delta=+0.044$, positive in 40 of 46 sections, $p<0.001$) and
in D3 (12 of 12 sections, $p=0.0005$); D4, D6 and D2 are significantly lower, and D1 and D7
show no effect. Restricted to precursor sections, D5 remains positive in 17 of 20, so the
signal is present before invasion. This is consistent with AT2 as the cell of origin in this
disease and with an independent supervised classifier whose signal was likewise confined to
AT2.

The score is a continuous quantity and not a malignancy call; no malignant label is assigned
to any spot or domain. Within D5, splitting spots into high and low copy-number groups
produces spatially coherent patches, but the split is not separable from RNA content: the
correlation between score and total UMI within D5 is $r=0.740$, and splitting by spot counts
alone reproduces roughly 75\% of the difference in AT2 content between the two groups. A
``D5-high'' subdomain therefore cannot be described as a tumour subdomain.

\subsection{Prognostic association of the niche domains}

The domain marker signatures were tested for prognostic information in an independent cohort
(TCGA-LUAD; 483 patients, 177 deaths; 381 patients and 143 events for disease-free survival).
Using each domain's top-30 markers standardised across samples, the associations are
consistent between a continuous parameterisation and a median split (Fig.~\ref{fig:2}f) and
between overall and disease-free survival. A high D7 (lymphoid) signature is associated with
better outcome (continuous HR 0.77, $p<0.001$; median split HR 0.60) and a high D3 signature
with worse (HR 1.46, $p<0.001$; HR 2.03); D4 and D5 are protective on overall survival.

The D3 result requires qualification. Its signature is drawn from the raw marker table,
which the depth analysis shows to be dominated by depth for this domain. Replacing it with
the depth-balanced signature reduces the adjusted hazard ratio from 2.30 to 1.30
($p=0.09$), and the connectivity-map signal for the same signature falls within the random
range. The direction and its robustness are therefore reported, but not the point estimate,
and D3 is not treated as an established prognostic feature.

\subsection{\emph{In-silico} reversal of the disease state}

To ask whether the disease state can be reversed, a query signature was derived from the
single-cell data and used against two perturbation libraries.

\paragraph{Query construction.} For each patient the difference in mean expression between
invasive and precursor cells was computed and the median taken across patients, giving a
16,104-gene axis. This axis was filtered against an independent survival signal (per-gene
adjusted Cox models in TCGA-LUAD), retaining a gene only where the direction of disease
change agreed with the direction of prognostic effect, so that a gene rising in disease is
retained only if higher expression also predicts worse outcome, and conversely
(Fig.~\ref{fig:3}a). The filter removes 9,270 of 18,069 genes (51\% of the full set; 44\% of
the 15,792 genes with survival data), of which 6,993 change in disease but are protective
and 2,277 lack survival annotation. It is applied because a reversal query that ignores it
would target protective programmes as readily as harmful ones.

\paragraph{Scoring schemes.} The signature was scored against the LINCS library using the
CMap weighted-KS statistic ($\tau$) and two correlation-based scores. The library was rebuilt
and held locally rather than queried through the CLUE scoring service, which was retired on
2026-01-31 (Methods), so the version queried is fixed and re-runnable and every scoring path
is cross-validated against the official implementation. The two correlation scores are
nearly redundant (Spearman $\rho=0.950$; 433 of 1,826 compounds shared in their top 500),
whereas the CMap statistic shares only about a third of its ranking with them
($\rho=0.36$--$0.38$; 18--19 of 50 in the top 50) (Fig.~\ref{fig:3}h). This reflects the
statistics rather than the implementation: CMap-type scores use only the membership of the
query's gene sets and discard the perturbation values, and are therefore systematically
different from value-based correlations. Requiring a compound to rank in the top 100 under
all three schemes is consequently a stricter filter than it appears, and yields
\textbf{34 compounds} (chance expectation 5.5), of which 16 (47\%) target PI3K/mTOR, with
smaller groups against HSP90, the proteasome and cell-cycle kinases (Fig.~\ref{fig:3}d; Table~\ref{tab:compounds}).

\begin{table*}[!ht]
\centering
\begin{tabular*}{\textwidth}{@{\extracolsep{\fill}}lrS[table-format=-1.3]S[table-format=-1.3]S[table-format=-1.3]rr@{}}
\toprule
\multicolumn{1}{@{}l}{\tabheadfont Compound} & \multicolumn{1}{l}{\tabheadfont Mechanism} & \multicolumn{1}{c}{\tabheadfont CMap $\tau$} & \multicolumn{1}{c}{\tabheadfont Cor-Sp $\rho$} & \multicolumn{1}{c}{\tabheadfont Cor-Pr $\rho$} & \multicolumn{1}{c}{\tabheadfont Rank CMap} & \multicolumn{1}{c@{}}{\tabheadfont Rank Cor-Sp} \\
\midrule
AZD-2014 & PI3K / mTOR & -0.355 & -0.186 & -0.112 & 1 & 15 \\
lacidipine & PI3K / mTOR & -0.330 & -0.188 & -0.113 & 2 & 14 \\
AZD-8055 & PI3K / mTOR & -0.329 & -0.176 & -0.111 & 3 & 23 \\
PKI-179 & PI3K / mTOR & -0.322 & -0.223 & -0.152 & 5 & 2 \\
GDC-0941 & PI3K / mTOR & -0.311 & -0.111 & -0.071 & 7 & 61 \\
MLN-0128 & PI3K / mTOR & -0.311 & -0.205 & -0.128 & 8 & 8 \\
NVP-AUY922 & HSP90 & -0.304 & -0.135 & -0.091 & 9 & 46 \\
GDC-0980 & PI3K / mTOR & -0.302 & -0.215 & -0.137 & 10 & 5 \\
BIIB-021 & JAK2 & -0.294 & -0.170 & -0.118 & 13 & 28 \\
taselisib & PI3K / mTOR & -0.291 & -0.146 & -0.116 & 14 & 41 \\
okadaic acid & phosphatase & -0.286 & -0.179 & -0.142 & 18 & 20 \\
XL-888 & cell-cycle kinase & -0.280 & -0.223 & -0.134 & 20 & 3 \\
JW-7-24-1 & mTOR & -0.279 & -0.094 & -0.065 & 22 & 87 \\
torin-2 & PI3K / mTOR & -0.277 & -0.237 & -0.153 & 23 & 1 \\
midostaurin & kinase & -0.276 & -0.109 & -0.059 & 24 & 62 \\
NVP-BEZ235 & PI3K / mTOR & -0.274 & -0.199 & -0.124 & 26 & 10 \\
bortezomib & proteasome & -0.271 & -0.180 & -0.106 & 29 & 19 \\
pentobarbital & GABA-A & -0.266 & -0.163 & -0.093 & 30 & 33 \\
PI-103 & PI3K / mTOR & -0.264 & -0.155 & -0.093 & 31 & 36 \\
TAK-285 & HER2 / EGFR & -0.261 & -0.117 & -0.081 & 32 & 57 \\
voxtalisib & PI3K / mTOR & -0.257 & -0.093 & -0.063 & 35 & 91 \\
torin-1 & PI3K / mTOR & -0.250 & -0.149 & -0.094 & 40 & 39 \\
PF-03814735 & Aurora kinase & -0.242 & -0.119 & -0.067 & 42 & 55 \\
AT-13387 & HSP90 & -0.237 & -0.154 & -0.089 & 48 & 37 \\
OSI-027 & mTOR & -0.236 & -0.093 & -0.062 & 51 & 89 \\
WYE-125132 & mTOR & -0.217 & -0.172 & -0.107 & 69 & 25 \\
MG-132 & proteasome & -0.214 & -0.184 & -0.109 & 72 & 16 \\
NVP-TAE226 & FAK & -0.213 & -0.164 & -0.101 & 73 & 32 \\
GSK-2110183 & PI3K / mTOR & -0.209 & -0.101 & -0.081 & 77 & 79 \\
dinaciclib & CDK & -0.208 & -0.177 & -0.102 & 82 & 22 \\
GSK-2126458 & PI3K / mTOR & -0.208 & -0.218 & -0.139 & 83 & 4 \\
YM-155 & survivin & -0.205 & -0.140 & -0.075 & 87 & 42 \\
tanespimycin & HSP90 & -0.205 & -0.183 & -0.118 & 88 & 17 \\
SB-939 & cell-cycle kinase & -0.199 & -0.102 & -0.059 & 94 & 77 \\
\bottomrule
\end{tabular*}
\caption{\textbf{The 34 compounds common to all three scoring schemes.} Compounds rank in
the top 100 under each of the CMap weighted-KS statistic ($\tau$) and the two correlation
scores, applied to the prognosis-filtered patient-internal disease axis against the LINCS
library (1{,}826 compounds). Rows are ordered by CMap rank. Values are medians across the
cell lines in which a compound was profiled; more negative indicates stronger reversal. The
chance expectation for a compound to appear in all three top-100 lists is 5.5. Mechanisms
are annotated from library metadata, curated for incompletely annotated compounds; the
assignment plays no part in selection.}
\label{tab:compounds}
\end{table*}


\paragraph{Reproducibility and modality independence.} Splitting patients at random into
halves, deriving the axis independently in each half and scoring both against LINCS
reproduces the ranking with a median Spearman $\rho$ of 0.962 across four splits, against a
permutation null of $\pm0.045$ (Fig.~\ref{fig:3}e, left); the top-500 lists share 438
compounds on average against a chance expectation of 137. One split is weaker ($\rho=0.752$)
and all four are reported. A second query axis was derived from the spatial data, after
depth balancing, by comparing invasive and precursor spots within nUMI deciles, and scored
against the same library with the same statistic. The two axes agree at Spearman
$\rho=0.735$ (top-100 overlap 52, chance expectation 5.5) (Fig.~\ref{fig:3}e, right). Since
the two axes derive from different platforms and different pairing strategies, this is a
stronger reproducibility check than repeating the scoring scheme within one modality.

\paragraph{Positive control.} Before the candidate list was examined, the pipeline was
verified against an official benchmark with a known answer, the epiblast-to-nascent-mesoderm
transition, using the authors' own implementation, manifold and standard deviations
\citep{pan2026scmg}. The known regulators are recovered at the top of the ranking: TBXT 1st
of 8,450, MSGN1 2nd, MIXL1 3rd, POU5F1 7th, SNAI1 9th, EOMES 11th, NANOG 12th and EVX1 15th
(Fig.~\ref{fig:3}b). This establishes that an empty result from the pipeline indicates the
query or the data rather than a broken method. Numerical identity with the official
implementations was confirmed separately: the weighted-KS implementation reproduces the
reference to a maximum absolute difference of $6.6\times10^{-14}$ with 2000 of 2000 signs
agreeing, and the vectorised causal score reproduces the authors' Python at Spearman
1.00000000. Table~\ref{tab:validation} collects these and the other numerical agreements
reported below.

\paragraph{The reversal acts on specific domains.} Using each domain's own depth-balanced
signature as a query, 26 of the 34 candidates fall in the top 100 for D3 and 21 for D5,
against a chance expectation of 1.9; the remaining five domains are at or near chance
(Fig.~\ref{fig:3}g). The compounds therefore act on the ECM/interstitial and AT2 domains
rather than on the tissue as a whole.

\begin{figure*}[!ht]\centering
  \includegraphics[width=\textwidth]{Fig3}
  \caption{\textbf{\emph{In-silico} reversal of the disease state.}
(\textbf{a}) The query signature is filtered on independent survival evidence: disease
direction against adjusted hazard ratio, with genes whose directions disagree set to zero.
(\textbf{b}) Positive control --- the official epiblast-to-nascent-mesoderm benchmark,
with the known regulators labelled and ranked (TBXT 1st of 8,450).
(\textbf{c}) Spatial copy-number score per domain, each domain compared with the remaining
spots of the same section so that all spots in a comparison share one library.
(\textbf{d}) Compound-level reversal: CMap rank against Cor-Spearman rank for the 1,826
compounds, with the 34 compounds common to all three scoring schemes highlighted, and their
mechanism composition.
(\textbf{e}) The candidate ranking is not modality-specific: a single-cell-derived and a
spatially derived query axis agree at Spearman $\rho=0.735$.
(\textbf{f}) A domain's marker stability scales with its sequencing depth
($\rho=-0.89$; D3 is the sole outlier and the only invasive-restricted domain).
(\textbf{g}) Of the 34 candidates, 26 fall in the top 100 for the ECM/interstitial domain
and 21 for the AT2 domain, against a chance expectation of 1.9.
(\textbf{h}) One library, three scores: the two correlation scores are near-redundant
($\rho=0.950$) whereas the CMap statistic shares only about a third of its ranking with
them.}\label{fig:3}
\end{figure*}


\subsection{Gating of the candidate gene list}

The same disease axis was scored against a gene-perturbation library of 20,345 perturbations
\citep{pan2026scmg} using the authors' causal predictor. The gating framework applied to the
resulting ranking is described here because it converts an ungoverned list into a short one
and because it exposes a failure mode that a naive ranking conceals.

A gene-perturbation score is informative only if the effect depends on the cell type being
perturbed. Every perturbation was therefore scored on a K562 control axis, and those whose
effect was as strong there as on the target axis were removed. Of 9,878 scored entries, 20
were flagged as generalising, 48 were specific to the non-K562 reference, 1,898 were neutral
and 7,912 lacked K562 data and are reported as untested. Applying this gate together with
the survival-direction filter leaves 25 entries supported by all four scoring routes, 111 by
three and 7 by two.

To test whether the candidates act on a compartment rather than on the tissue as a whole,
the disease axis was rebuilt for each of 31 subtypes separately and every perturbation
re-ranked within each axis. Of 1,748 entries reaching the top 100 in at least one cell-type
axis, only two reach it in twenty or more of the 31. One is \emph{SFTPC} knockdown, which
ranks in the top 100 in 22 of 31 axes but fails the direction check, since lowering a
transcript whose high expression is protective does not describe a therapeutic direction.
The other is \emph{STAT1} knockdown (21 of 31 axes; hazard ratio 1.08, consistent with the
required direction). The remaining candidates are largely axis-specific. The headline number
of the ungated ranking is misleading for this reason: ubiquity across cell types is not
evidence of a target, and the entry that appears most broadly reproducible does not survive
the direction filter.

\subsection{Analyses that did not work}

Six further negative results are reported, each with the reason it is believed to have
failed.

\paragraph{Gene-level prediction on the pooled axis.} Applied to the tissue-pooled disease
axis, the pipeline produced a maximum causal score of 0.083 against 1.279 for the official
benchmark on the same implementation, and the distribution on our axis has no high tail,
whereas the benchmark's does. Since the method is demonstrably functional, the pooled axis
has no detectable match to the expression shifts recorded in the library. The
compartment-resolved query does rank, but its survivor count is one. A separate,
self-implemented cosine projection, which omits the causal score's gene-shift term, also
ranked genes, but its significance saturates at the permutation floor (237 of 2,439 genes at
the minimum attainable $p$) and it cannot discriminate.

\paragraph{A developmental trajectory that runs the wrong way.} Using the source study's
approach (monocle3 \citep{cao2019monocle3} with CytoTRACE \citep{gulati2020cytotrace}), the
score is highest rather than lowest in invasive disease: medians are Normal 0.288, AAH
0.374, AIS 0.288, MIA 0.277 and IAC 0.453, a non-monotonic profile with IAC highest. The
reversal is robust, with 22 of 23 patients showing a positive per-patient slope and a
depth-balanced audit finding IAC highest within every depth stratum. The measurement is
nevertheless ambiguous in a way that cannot be resolved. CytoTRACE is constructed from the
number of detected genes, so a high score in invasive disease is equally consistent with
greater transcriptional disorder and with a less differentiated state; two readings fit the
same numbers. Because the analysis measures a stage axis rather than a subtype axis, no
AIC/KAC or tumour labels being available, it can be interpreted neither as evidence against
the source study's trajectory nor as evidence for a different one.

\paragraph{A spectral mechanism that was an artefact of construction.} An initial attempt to
read the trajectory from the spectrum of an optimal-transport operator
\citep{schiebinger2019wot,lange2022cellrank} was abandoned after inspection. Under the stage
partition the coupling matrix is block upper triangular, so transport does not enter the
spectrum; in small independent checks with known row sums the eigenvalue union differed by
$0.000\times10^{0}$ and the number of near-zero eigenvalues equalled the cell count of the
first four stages. The apparent spectral gap was set by a coupling-weight parameter and not
by the data. The result was, on inspection, a property of how the matrix was assembled.

\paragraph{H\&E imaging at this resolution.} A pathology-language model
\citep{huang2023plip} was applied to the H\&E images to ask whether the morphological sequence
could be read directly from the histology. Across twelve independent configurations (six
feature arms crossed with two prompt sets) the best correlation was 0.729, and no setting
exceeded 0.74. The ceiling is set by the model architecture: PLIP takes a 224-pixel input
yielding 49 tokens, so one token spans about 181~$\mu$m of tissue, against the finest
available colour source at 5.670~$\mu$m per pixel. Above that scale the discriminative power
did not vary with detail, and finer sources cannot be tested in this cohort because none
exists.

For the use this study would make of it, the decisive result is that the model did not
resolve the precursor stages. Per-stage means for AAH, AIS and MIA all fell within 0.15 to
0.22 and interleaved with one another, so the three stages that define the sequence are not
separated by this model. What increased with image detail was only the contrast between
invasive disease and everything else, with the benign end depressed, rather than a gradient
across stages.

A search over 1,344 template-and-term combinations identifies what the model is in fact
measuring. The highest raw discrimination, an AUC of 0.972, came from the term
``invasive adenocarcinoma with desmoplastic stroma'', but that term also correlated most
strongly with tissue density ($\rho=0.574$), and projecting the density direction out of its
embedding reduced the AUC to 0.759. A weaker prompt chosen to be density-neutral retained
more of its discrimination under the same projection (raw AUC 0.833, $\rho=0.129$;
projected 0.791). Two independent measurements agree: the lesion score correlates at
$\rho=0.718$ with the tissue fraction of a spot and at $\rho=-0.794$ with an empty-slide
control; and the model does not merely fail to rank precursor lesions above normal lung, it
inverts the comparison, with AAH versus Normal AUCs of 0.28 to 0.32 across prompt sets.
Tumour tissue is itself denser than normal lung, so removing density also removes part of the
true signal and the projected values are a lower bound rather than a target. The relevant
point is that performance here is not separable from how much tissue a spot contains.

To establish that the ceiling reflects resolution rather than the choice of image, every
section was registered from the high-resolution image into the CytAssist frame (56 of 56
passing; Methods) and the same spots were re-scored from four image sources matched by
physical size. The finer source did not improve performance: slice-level correlation was
0.698 for the baseline image, 0.228 for CytAssist, 0.540 for CytAssist contrast-matched to
the baseline, and 0.215 for CytAssist downsampled to the baseline's pixel density. Because
downsampling and raw CytAssist perform alike, the difference is not one of pixel density but
of white balance and contrast between the two products, and matching the contrast recovers
most of it.

The conclusion is correspondingly narrow. At this image scale the model cannot perform the
task this study would require of it, which is to place a lesion on the precursor-to-invasive
axis. That is not a statement that histological detail is uninformative in general: testing
that would require a finer colour source than this cohort provides.

\paragraph{The ligand layer.} NicheNet was used with fibroblasts as receivers, the ten most
abundant non-fibroblast cell types of the ECM/interstitial domain as senders, and the same
data-derived target set. Only five ligands passed the expression and network gates, and none
produced a ligand--target link in the output matrix. Their ligand activities were
correspondingly flat, the best corrected area under the precision--recall curve being 0.093
against a baseline of zero. Since NicheNet's own baseline correction was applied, the
failure lies in the candidate set rather than in the scoring: with five ligands the sender
panel is too thin for the network to resolve anything. Ligand-level inference is the natural
next step, and this is the specific reason it was not pursued here.

\paragraph{A finer spatial resolution that measured sections.} One attempt to ask whether the
RCTD composition space supports more than seven domains pooled spots from all 56 sections
into a single $k$-means and cross-validated the result. It returned a clean optimum, but the
clusters tracked section identity rather than niche. Pooling across sections makes
between-section differences the dominant axis of variation, so the bootstrap measured
whether clusters mixed across sections are stable and not whether the composition space
supports a finer partition. The analysis was superseded by a per-section clustering with the
same $K$ scan, which is the version reported above.

\begin{table*}[!ht]
\centering
\begin{tabularx}{\textwidth}{@{}>{\itshape\raggedright\arraybackslash}p{94pt}>{\raggedright\arraybackslash}p{152pt}X@{}}
\toprule
\tabheadfont Analysis & \tabheadfont What was observed & \tabheadfont Most likely reason, and the check behind it \\
\midrule
Single-cell copy number & no usable malignant/non-malignant separator across six tool families; the one AUROC of 0.5000 was a degeneracy (a single score value for all 19{,}748 cells) & the per-patient benign anchor is unstable: the same 3{,}890 normal cells gave aneuploid calls of 0\%, 13.6\% and 49.8\% in three runs. Not evidence that copy-number differences are absent. \\
\addlinespace[2pt]
Cohort-level spatial copy number & the invasive-vs-precursor difference shrank 71\% after depth matching ($p = 0.092$, CI including zero) & depth and tissue density are the same measurement here. A random subsample reproduced the pre-matching result index by index, so the original effect was depth-driven rather than a sample-size artefact. \\
\addlinespace[2pt]
Copy-number subdomain within AT2 & the score correlates with total UMI at $r = 0.755$, and a count-only split reproduces $\sim$75\% of the AT2 contrast & RNA content and copy-number score are not separable; a high-score subdomain cannot be called a tumour subdomain. \\
\addlinespace[2pt]
Developmental trajectory & the score is highest, not lowest, in invasive disease (0.453 vs 0.288--0.374), robust in 22 of 23 patients & we measure a stage axis, not a subtype axis (no AIC/KAC or tumour labels). The score is built from detected-gene counts, so the same numbers also fit ``more transcriptionally disordered''. \\
\addlinespace[2pt]
Transport-spectrum trajectory & the apparent spectral gap was set by a coupling-weight parameter & under a stage partition the coupling matrix is block upper triangular, so transport never enters the spectrum. A parameter was being read as biology. \\
\addlinespace[2pt]
Gene-level reversal, pooled axis & maximum causal score 0.083 against 1.279 for the official benchmark & the pooled axis has no detectable match in the library. The positive control passes, so this is a statement about the query. \\
\addlinespace[2pt]
Gene-level reversal, per compartment & only two perturbations reach the top 100 in $\geq$20 of 31 subtype axes, and one fails the direction check & ubiquity across compartments is not evidence of a target; the survivor list is a single entry and is reported as an observation. \\
\addlinespace[2pt]
Ligand layer & five ligands passed the expression and network gates; none produced a ligand--target link; best corrected AUPR 0.093 & the sender panel is too thin. The scoring is baseline-corrected, so this is not a scoring artefact. \\
\addlinespace[2pt]
Four spatial programmes by stage & after depth matching the ordering is identical across all ten deciles but does not follow disease progression & depth creates a spurious trend. Reported as a negative control, not as a result. \\
\addlinespace[2pt]
H\&E imaging & best correlation 0.729 across twelve configurations; the three precursor stages fall within 0.15--0.22 and interleave (AAH versus Normal AUC 0.28--0.32); the top-scoring prompt correlates with tissue density at $\rho=0.574$ and falls to 0.759 when that direction is projected out & architectural resolution (181\,$\mu$m per token against 5.670\,$\mu$m per pixel). At this scale the model's performance is not separable from tissue density and it cannot place a lesion on the precursor-to-invasive axis. Does not mean histological detail is uninformative in general. \\
\addlinespace[2pt]
Domain stability gate & cross-seed ARI never reached 0.90 (median 0.782) & the domain definition is method-sensitive (BANKSY 7, RCTD 6, marker-based 5; ARI 0.352). The partition is reported as exploratory. \\
\addlinespace[2pt]
Pooled cross-section clustering & a clean optimum that tracked section identity & pooling sections makes between-section variation the dominant axis; superseded by per-section clustering. \\
\bottomrule
\end{tabularx}
\caption{\textbf{Negative results and the reason we believe each one failed.} The right-hand
column is an assessment, not a measurement: it names the most likely explanation and, where
we could test it, the check that separates an absent signal from a broken method. The
distinction matters because a negative result with a passing positive control constrains the
biology, whereas one without does not.}
\label{tab:negatives}
\end{table*}


\begin{table*}[!ht]
\centering
\begin{tabularx}{\textwidth}{@{}>{\raggedright\arraybackslash}p{116pt}>{\raggedright\arraybackslash}p{104pt}X@{}}
\toprule
\tabheadfont Check & \tabheadfont Compared against & \tabheadfont Result \\
\midrule
\tabgroup{3}{Numerical agreement with the reference implementation}
Weighted-KS statistic & \texttt{signatureSearch} 1.12.0 & max.\ abs.\ difference \num{6.6e-14}; signs agree \num{2000}/\num{2000} \\
Normalised connectivity & \texttt{signatureSearch} 1.12.0 & max.\ abs.\ difference \num{1.2e-13} \\
Causal gene score & authors' Python & Spearman \num{1.00000000}; signs agree 100\% (\num{19819} entries) \\
\addlinespace[3pt]
\tabgroup{3}{Biological positive control}
Official benchmark & epiblast $\rightarrow$ nascent mesoderm & TBXT 1st of \num{8450}; MSGN1 2nd; MIXL1 3rd; POU5F1 7th; SNAI1 9th; EOMES 11th; NANOG 12th; EVX1 15th \\
\addlinespace[3pt]
\tabgroup{3}{Stability and agreement}
Split-half & 4 patient partitions & median $\rho = 0.962$ (one split 0.752); permutation null $\pm0.045$ \\
Cross-modality & single-cell axis vs spatial axis & $\rho = 0.735$; top-100 overlap 52 (chance 5.5) \\
Scoring schemes & CMap vs the two correlations & Cor--Cor $\rho = 0.950$ (433 of top 500 shared); CMap--Cor $\rho = 0.36$--$0.38$ \\
\addlinespace[3pt]
\tabgroup{3}{Geometric self-consistency}
Scale-bar calibration & adjacent-spot spacing $= 100$\,$\mu$m & 1\,px $= 0.2513$\,$\mu$m; full panel 10.9\,mm (11\,mm format) \\
\bottomrule
\end{tabularx}
\caption{\textbf{Method verification and positive controls.} Every numerical check compares
our own implementation with the reference implementation on identical input, so a
discrepancy would be an implementation error rather than a modelling choice. The biological
control establishes that an empty result from this pipeline reflects the query and not the
method.}
\label{tab:validation}
\end{table*}

